\section{Discussion}

\subsection{Coverage breadth and facility scale}
The central interpretive distinction is between how widely a network is observed and how much of its facility stock it contains. In the 2022--2023 endpoint frame, 117 of 225 cities are $00$, yet they account for only 3.47\% of endpoint facilities; the 93 $++$ cities account for 95.04\%. City counts therefore describe geographic reach within the archive, whereas facility shares describe the scale carried by the observed positive network. Because the 117 cities are numerous but facility-light, coverage is a geographic-reach objective; throughput should be assessed from the activity total alongside facility-normalized intensity, both of which are dominated by the facility-bearing positive network. The two statistics should be reported together: neither converts an operator-network frame into a national coverage estimate, and neither treats a source-coded zero as proof of absent demand or an uncommissioned site.

This separation extends established infrastructure and adoption work. Infrastructure and adoption can reinforce one another through network effects, with the relationship depending on subsidy design and local context \citep{li2017,springel2021,burra2024,hu2025_charging}. Conversely, station studies show that utilization, occupancy, delivered energy and profitability depend on charger type, location and operating conditions \citep{borlaug2023,hecht2022,zhao2024_utilization,du2024_accessibility,zheng2025_landuse}. A city-level recorded-activity ratio consequently occupies a different evidential position from vehicle adoption, access, energy per port or a cost-based utilization measure. The present network result refines these literatures by showing how the conclusion changes when the unit of interpretation is city reach, observed facility scale or facility-weighted network intensity.

The same point applies to zeros and positive records. Positive-record frequency becomes more common in the long cohort even while facilities grow faster than recorded activity, and the broad endpoint decomposition assigns most of the change to persistently positive records. This pattern is consistent with a network that expands its observed footprint while the relationship between stock and recorded use varies across locations. It does not decide whether low activity reflects demand, access, reporting, equipment mix or other local conditions. Studies of censored demand, habitual use and reliability likewise caution that delivered activity, charging opportunity and service difficulty are distinct outcomes \citep{huttel2023,weekx2024,weekx2025_habits,patil2023,yu2025_reliability}. For policy evaluation, a coverage target and a throughput target therefore require different denominators and different evidence.

\subsection{Growth, proportionality and persistence}
The fixed-city evidence separates three questions that are often collapsed into one ratio. First, most long-cohort cities record activity growth while the facility indicator grows faster; second, a smaller group combines facility expansion with absolute activity contraction; third, the relative-growth category can change within the same city. In the primary pairing, 17 of 27 early faster-growth cities change category by the later window. These paired counts make persistence observable and show why a single window label cannot stand in for sustained performance: evaluation should track both the relative pace of activity and facility growth and whether that pace persists across a prespecified subsequent window, without treating a label as a development stage or assigning a cause to the change.

The annual and quarterly models provide a related but distinct scale check. The two annual PPML estimates are positive and below unit elasticity on their respective supports, whereas the log-OLS estimate is \AnnualOLSBeta\ [\AnnualOLSLow, \AnnualOLSHigh] and its unit-elasticity test gives $p=\AnnualOLSPOne$, so the log-scale interval does not exclude proportionality. The three-quarter endpoint is likewise subproportional. The $c_3-c_0$ interval spans zero, so the data resolve a finite-horizon proportionality gap more clearly than an additional delayed increment. This complements station-level evidence that activity depends on equipment, built environment, temporal routines and accessibility \citep{wolbertus2018,zhao2024_utilization,du2024_accessibility,zheng2025_landuse}; it does not identify which of those factors generates a particular city trajectory.

The contribution and recorded-week decompositions clarify why different summaries can coexist. A city can improve its own intensity yet contribute negatively to facility-weighted network intensity when its facility weight changes; equal-city means can be positive while the changes in median and facility-weighted network intensity are negative. In the fixed 52-city check, later negative contributions are spread across more cities and are less concentrated by the stated HHI and top-share measures. These are accounting results, but they make the evaluation choice explicit: geographic reach, network activity totals, facility-weighted intensity and persistence are distinct measures. A programme that supports coverage in low-throughput places should not be judged by the same denominator as a programme whose stated objective is high network utilization.

The proportionality and timing results are conditional benchmarks, not forecasts of a city-wide demand response: annual estimates condition on city and year effects, weather and observed exposure, whereas the quarterly estimate is a finite endpoint under a stipulated persistent facility difference. Read as evaluation targets, the results distinguish coverage (where facilities are observed), activity growth (whether recorded use keeps pace with the observed stock) and persistence (whether a city's relative pace holds over a prespecified horizon). The three-quarter endpoint is therefore a horizon-specific benchmark, not a universal reporting deadline or maturity schedule; category changes do not establish a common recovery path, and the pandemic-context comparison remains a conditioning-index comparison rather than a causal pandemic effect.

\subsection{Public support and operating responsibilities}
The policy text shows why construction support and operating evaluation should not be represented by one treatment count. The policy-mix perspective emphasizes the interaction of instruments and their assigned functions \citep{rogge2016}; Chinese governance studies likewise distinguish subsidy design, operational performance and stakeholder responsibilities \citep{hu2026,zhang2024}. In the present corpus, information and supervision are more widespread than operation and maintenance or continuing-operation/payment clauses, while the functions overlap. This pattern describes administrative emphasis, not implementation intensity.

The functional distinction translates this emphasis into verifiable operating responsibilities. An information requirement can be paired with evidence of platform completeness and reporting continuity; a maintenance requirement with equipment-condition and response-time records; and a continuing-operation or payment condition with dated service evidence and an eligibility period. These pairings inspect the administrative objects named in the documents, rather than treating their frequency as implementation. They accord with policy analyses that treat public charging as a system of coordination, data sharing and enforceable management arrangements \citep{tilly2025}. Coverage, activity and continuity can then be assessed using reach measures, compatible activity records and interruption or operating records, respectively.

The quantitative policy results reinforce that separation. The information/supervision association with recorded activity is compatible with greater reporting visibility, accompanying service change or both; its positive coefficient does not establish an information-mediated activity effect. Operation/maintenance intervals include zero and vary in width across outcomes. The quarterly extension has a facility-lead diagnostic showing that pre-policy trajectories are not comparable to the benchmark; the diagnostic is informative about identification, but its lead and post-policy contrast use different reference structures and should not be subtracted or connected as a causal effect. Together, the text, annual association and timing extension support differentiated verification: check whether the prescribed information is available, whether equipment and access are maintained, and whether continuity or payment conditions are met over the stated period. The appropriate evaluation criterion follows from the specified responsibility and its observation period, rather than a universal activity ratio or operating deadline.
